
\documentclass{article}
\usepackage{colortbl}
\usepackage{makecell}
\usepackage{multirow}
\usepackage{supertabular}

\begin{document}

\newcounter{utterance}

\centering \large Interaction Transcript for game `function\_detective', experiment `two\_numbers\_passive\_labeled\_pairs', episode 249 with Qwen3.6{-}35B{-}A3B{-}FP8{-}without{-}reasoning.
\vspace{24pt}

{ \footnotesize  \setcounter{utterance}{1}
\setlength{\tabcolsep}{0pt}
\begin{supertabular}{c@{$\;$}|p{.15\linewidth}@{}p{.15\linewidth}p{.15\linewidth}p{.15\linewidth}p{.15\linewidth}p{.15\linewidth}}
    \# & \multicolumn{2}{c}{Player} && \multicolumn{2}{c}{Game Master} \\
    \hline

    \theutterance \stepcounter{utterance}  
    & & & \multicolumn{4}{p{0.6\linewidth}}{
        \cellcolor[rgb]{0.9,0.9,0.9}{
            \makecell[{{p{\linewidth}}}]{
                \texttt{\tiny{[P1$\langle$GM]}}
                \texttt{I have a hidden Python function for you to figure out.} \\
\\ 
\texttt{Your job is to infer what it does and then submit a correct Python implementation.} \\
\\ 
\texttt{Function signature:} \\
\texttt{(a: int, b: int) {-}> int} \\
\\ 
\texttt{You have 15 turns total.} \\
\\ 
\texttt{I reveal candidate input{-}output pairs one at a time and tell you whether each one is correct.} \\
\\ 
\texttt{You must make your best SOLVE attempt no later than the second{-}to{-}last turn.} \\
\texttt{Do not wait until the very end.} \\
\\ 
\texttt{You can use these actions:} \\
\texttt{{-} JSON format described at the end of the prompt} \\
\texttt{{-} SOLVE: ```python ... ```} \\
\\ 
\texttt{Do not add import statements in your guess.} \\
\texttt{Anything you might need is already available, so just write the function body in hypothesis and return the answer.} \\
\\ 
\texttt{Preloaded examples:} \\
\texttt{1. ((({-}10, 10), 0), True)} \\
\\ 
\texttt{Scoring:} \\
\texttt{{-} correctness is required for any non{-}zero score} \\
\texttt{{-} in interactive modes, solving with fewer information requests is better} \\
\\ 
\texttt{Input formatting reminders:} \\
\texttt{{-} strings must be quoted} \\
\texttt{{-} booleans must be True/False} \\
\texttt{{-} use commas between multiple values} \\
\texttt{{-} do not write any extra action tags in the same message} \\
\\ 
\texttt{Format example:} \\
\texttt{You write:} \\
\texttt{\{'current\_mode': <probing/confirming/don't know; choose whichever best describes your current state in the task>, 'input\_rationale': <brief description of why you chose this hypothesis>, 'hypothesis': <either the exact string "unknown" if you have no guess yet OR your current best guess as runnable Python matching the given signature, with a one{-}line docstring stating what it does. Never write placeholder code, if you cannot commit to a guess, use "unknown". Match the given signature.} \\
\texttt{e.g.} \\
\texttt{def f(x):} \\
\texttt{	"""Returns the absolute value of x."""} \\
\texttt{	return x if x >= 0 else {-}x} \\
\\ 
\texttt{Use "unknown" if you have no guess.>, 'confidence\_score': <score between 1{-}100. 100 means highest confidence in your hypothesis>\}} \\
\\ 
\texttt{Output only one valid json per turn.} \\
            }
        }
    }
     \\ \\

    \theutterance \stepcounter{utterance}  
    & \multicolumn{4}{p{0.6\linewidth}}{
        \cellcolor[rgb]{0.561,0.737,0.561}{
            \makecell[{{p{\linewidth}}}]{
                \texttt{\tiny{[P1$\rangle$GM]}}
                \texttt{\{'current\_mode': 'probing', 'input\_rationale': 'I have one preloaded example: f({-}10, 10) = True. This suggests the function might return a boolean, but the signature says it returns an int. In Python, True is equivalent to 1 and False to 0 when returned as an int. So f({-}10, 10) likely returns 1. I need more examples to understand the logic. Let me try some simple cases like f(0, 0), f(1, 0), f(0, 1), f(1, 1) to see if it relates to comparison, arithmetic, or bitwise operations.', 'hypothesis': 'unknown', 'confidence\_score': 10\}} \\
            }
        }
    }
    & & \\ \\

    \theutterance \stepcounter{utterance}  
    & & & \multicolumn{4}{p{0.6\linewidth}}{
        \cellcolor[rgb]{0.9,0.9,0.9}{
            \makecell[{{p{\linewidth}}}]{
                \texttt{\tiny{[P1$\langle$GM]}}
                \texttt{Game Violation: ParseError: Response is not valid JSON.} \\
\texttt{Output ONLY:} \\
\texttt{One JSON with confidence\_score, current\_mode, input\_rationale, and hypothesis keys ...} \\
\texttt{OR} \\
\texttt{SOLVE: ```python ... ```} \\
            }
        }
    }
     \\ \\

    \theutterance \stepcounter{utterance}  
    & \multicolumn{4}{p{0.6\linewidth}}{
        \cellcolor[rgb]{0.561,0.737,0.561}{
            \makecell[{{p{\linewidth}}}]{
                \texttt{\tiny{[P1$\rangle$GM]}}
                \texttt{\{} \\
\texttt{  "current\_mode": "probing",} \\
\texttt{  "input\_rationale": "I have one preloaded example: f({-}10, 10) = True. Since the return type is int, True likely maps to 1. I need to probe more to understand the relationship between a and b. Let's try f(0, 0) to see if it returns 0 or 1, and f(1, 0) to check asymmetry.",} \\
\texttt{  "hypothesis": "unknown",} \\
\texttt{  "confidence\_score": 10} \\
\texttt{\}} \\
            }
        }
    }
    & & \\ \\

    \theutterance \stepcounter{utterance}  
    & & & \multicolumn{4}{p{0.6\linewidth}}{
        \cellcolor[rgb]{0.9,0.9,0.9}{
            \makecell[{{p{\linewidth}}}]{
                \texttt{\tiny{[P1$\langle$GM]}}
                \texttt{OUTPUT: (((2, {-}1), {-}1), False)} \\
            }
        }
    }
     \\ \\

    \theutterance \stepcounter{utterance}  
    & \multicolumn{4}{p{0.6\linewidth}}{
        \cellcolor[rgb]{0.561,0.737,0.561}{
            \makecell[{{p{\linewidth}}}]{
                \texttt{\tiny{[P1$\rangle$GM]}}
                \texttt{\{} \\
\texttt{  "current\_mode": "probing",} \\
\texttt{  "input\_rationale": "I have two examples now: f({-}10, 10) = True (1) and f(2, {-}1) = False (0). Let's analyze: {-}10 < 10 is True, 2 > {-}1 is True. So it's not simply a < b. Let's try f(0, 0) to see if equal inputs return 0 or 1, and f(1, 2) to check if a < b returns 1.",} \\
\texttt{  "hypothesis": "unknown",} \\
\texttt{  "confidence\_score": 10} \\
\texttt{\}} \\
            }
        }
    }
    & & \\ \\

    \theutterance \stepcounter{utterance}  
    & & & \multicolumn{4}{p{0.6\linewidth}}{
        \cellcolor[rgb]{0.9,0.9,0.9}{
            \makecell[{{p{\linewidth}}}]{
                \texttt{\tiny{[P1$\langle$GM]}}
                \texttt{OUTPUT: ((({-}10, 1000), {-}10000), True)} \\
            }
        }
    }
     \\ \\

    \theutterance \stepcounter{utterance}  
    & \multicolumn{4}{p{0.6\linewidth}}{
        \cellcolor[rgb]{0.561,0.737,0.561}{
            \makecell[{{p{\linewidth}}}]{
                \texttt{\tiny{[P1$\rangle$GM]}}
                \texttt{\{} \\
\texttt{  "current\_mode": "probing",} \\
\texttt{  "input\_rationale": "I have three examples: f({-}10, 10) = True, f(2, {-}1) = False, f({-}10, 1000) = True. Let's look at the relationship between a and b. In the first case, {-}10 and 10: maybe it's checking if a + b == 0? {-}10 + 10 = 0. In the second case, 2 + ({-}1) = 1 != 0. In the third case, {-}10 + 1000 = 990 != 0. So it's not a + b == 0. Let's try f(0, 0) to see if 0 + 0 == 0 returns True. Also, let's try f(5, {-}5) to check if a + b == 0.",} \\
\texttt{  "hypothesis": "unknown",} \\
\texttt{  "confidence\_score": 10} \\
\texttt{\}} \\
            }
        }
    }
    & & \\ \\

    \theutterance \stepcounter{utterance}  
    & & & \multicolumn{4}{p{0.6\linewidth}}{
        \cellcolor[rgb]{0.9,0.9,0.9}{
            \makecell[{{p{\linewidth}}}]{
                \texttt{\tiny{[P1$\langle$GM]}}
                \texttt{OUTPUT: ((({-}2, {-}99), 199), False)} \\
            }
        }
    }
     \\ \\

    \theutterance \stepcounter{utterance}  
    & \multicolumn{4}{p{0.6\linewidth}}{
        \cellcolor[rgb]{0.561,0.737,0.561}{
            \makecell[{{p{\linewidth}}}]{
                \texttt{\tiny{[P1$\rangle$GM]}}
                \texttt{\{} \\
\texttt{  "current\_mode": "probing",} \\
\texttt{  "input\_rationale": "I have four examples: f({-}10, 10) = True, f(2, {-}1) = False, f({-}10, 1000) = True, f({-}2, {-}99) = False. Let's look for a pattern. Maybe it's checking if a and b have the same sign? {-}10 and 10 have different signs {-}> True. 2 and {-}1 have different signs {-}> False. This doesn't hold. Maybe it's checking if a * b < 0? {-}10 * 10 = {-}100 < 0 {-}> True. 2 * {-}1 = {-}2 < 0 {-}> False. This doesn't hold either. Let's try f(0, 0) to see if zero is involved. Also, let's try f(5, 5) to check if same sign/positive returns False.",} \\
\texttt{  "hypothesis": "unknown",} \\
\texttt{  "confidence\_score": 10} \\
\texttt{\}} \\
            }
        }
    }
    & & \\ \\

    \theutterance \stepcounter{utterance}  
    & & & \multicolumn{4}{p{0.6\linewidth}}{
        \cellcolor[rgb]{0.9,0.9,0.9}{
            \makecell[{{p{\linewidth}}}]{
                \texttt{\tiny{[P1$\langle$GM]}}
                \texttt{OUTPUT: ((({-}2, {-}10), 20), True)} \\
            }
        }
    }
     \\ \\

    \theutterance \stepcounter{utterance}  
    & \multicolumn{4}{p{0.6\linewidth}}{
        \cellcolor[rgb]{0.561,0.737,0.561}{
            \makecell[{{p{\linewidth}}}]{
                \texttt{\tiny{[P1$\rangle$GM]}}
                \texttt{\{} \\
\texttt{  "current\_mode": "probing",} \\
\texttt{  "input\_rationale": "I have five examples now: f({-}10, 10) = True, f(2, {-}1) = False, f({-}10, 1000) = True, f({-}2, {-}99) = False, f({-}2, {-}10) = True. Let's look at the relationship between a and b again. Maybe it's checking if a divides b or b divides a? {-}10 divides 10? Yes. 2 divides {-}1? No. {-}10 divides 1000? Yes. {-}2 divides {-}99? No. {-}2 divides {-}10? Yes. This looks promising! Let's verify: a | b (a divides b) or b | a? In all True cases, at least one divides the other. In False cases, neither divides the other. Let's test f(3, 6) which should be True, and f(3, 4) which should be False.",} \\
\texttt{  "hypothesis": "unknown",} \\
\texttt{  "confidence\_score": 40} \\
\texttt{\}} \\
            }
        }
    }
    & & \\ \\

    \theutterance \stepcounter{utterance}  
    & & & \multicolumn{4}{p{0.6\linewidth}}{
        \cellcolor[rgb]{0.9,0.9,0.9}{
            \makecell[{{p{\linewidth}}}]{
                \texttt{\tiny{[P1$\langle$GM]}}
                \texttt{OUTPUT: (((1000, 0), {-}1), False)} \\
            }
        }
    }
     \\ \\

    \theutterance \stepcounter{utterance}  
    & \multicolumn{4}{p{0.6\linewidth}}{
        \cellcolor[rgb]{0.561,0.737,0.561}{
            \makecell[{{p{\linewidth}}}]{
                \texttt{\tiny{[P1$\rangle$GM]}}
                \texttt{\{} \\
\texttt{  "current\_mode": "probing",} \\
\texttt{  "input\_rationale": "I have six examples. The divisibility hypothesis (a divides b or b divides a) seems strong but let's check the new case: f(1000, 0) = False. 1000 divides 0? Yes, 0 = 1000 * 0. So if the hypothesis was 'a divides b', this should be True. But it's False. What about 'b divides a'? 0 divides 1000? No (division by zero). So if the condition is 'b divides a', then f(1000, 0) would be False. Let's re{-}check previous examples with 'b divides a': f({-}10, 10): 10 divides {-}10? Yes. f(2, {-}1): {-}1 divides 2? Yes. But f(2, {-}1) is False. So 'b divides a' is not it. Let's re{-}check 'a divides b': f({-}10, 10): {-}10 divides 10? Yes. f(2, {-}1): 2 divides {-}1? No. f({-}10, 1000): {-}10 divides 1000? Yes. f({-}2, {-}99): {-}2 divides {-}99? No.} \\
            }
        }
    }
    & & \\ \\

    \theutterance \stepcounter{utterance}  
    & & & \multicolumn{4}{p{0.6\linewidth}}{
        \cellcolor[rgb]{0.9,0.9,0.9}{
            \makecell[{{p{\linewidth}}}]{
                \texttt{\tiny{[P1$\langle$GM]}}
                \texttt{Game Violation: ParseError: Response is not valid JSON.} \\
\texttt{Output ONLY:} \\
\texttt{One JSON with confidence\_score, current\_mode, input\_rationale, and hypothesis keys ...} \\
\texttt{OR} \\
\texttt{SOLVE: ```python ... ```} \\
            }
        }
    }
     \\ \\

    \theutterance \stepcounter{utterance}  
    & \multicolumn{4}{p{0.6\linewidth}}{
        \cellcolor[rgb]{0.561,0.737,0.561}{
            \makecell[{{p{\linewidth}}}]{
                \texttt{\tiny{[P1$\rangle$GM]}}
                \texttt{\{} \\
\texttt{  "current\_mode": "probing",} \\
\texttt{  "input\_rationale": "I have six examples. The divisibility hypothesis (a divides b or b divides a) seems strong but let's check the new case: f(1000, 0) = False. 1000 divides 0? Yes, 0 = 1000 * 0. So if the hypothesis was 'a divides b', this should be True. But it's False. What about 'b divides a'? 0 divides 1000? No (division by zero). So if the condition is 'b divides a', then f(1000, 0) would be False. Let's re{-}check previous examples with 'b divides a': f({-}10, 10): 10 divides {-}10? Yes. f(2, {-}1): {-}1 divides 2? Yes. But f(2, {-}1) is False. So 'b divides a' is not it. Let's re{-}check 'a divides b': f({-}10, 10): {-}10 divides 10? Yes. f(2, {-}1): 2 divides {-}1? No. f({-}10, 1000): {-}10 divides 1000? Yes. f({-}2, {-}99): {-}2 divides {-}99? No.} \\
            }
        }
    }
    & & \\ \\

    \theutterance \stepcounter{utterance}  
    & & & \multicolumn{4}{p{0.6\linewidth}}{
        \cellcolor[rgb]{0.9,0.9,0.9}{
            \makecell[{{p{\linewidth}}}]{
                \texttt{\tiny{[P1$\langle$GM]}}
                \texttt{Game Violation: ParseError: Response is not valid JSON.} \\
\texttt{Output ONLY:} \\
\texttt{One JSON with confidence\_score, current\_mode, input\_rationale, and hypothesis keys ...} \\
\texttt{OR} \\
\texttt{SOLVE: ```python ... ```} \\
            }
        }
    }
     \\ \\

    \theutterance \stepcounter{utterance}  
    & \multicolumn{4}{p{0.6\linewidth}}{
        \cellcolor[rgb]{0.561,0.737,0.561}{
            \makecell[{{p{\linewidth}}}]{
                \texttt{\tiny{[P1$\rangle$GM]}}
                \texttt{\{} \\
\texttt{  "current\_mode": "probing",} \\
\texttt{  "input\_rationale": "I have six examples. The divisibility hypothesis (a divides b or b divides a) seems strong but let's check the new case: f(1000, 0) = False. 1000 divides 0? Yes, 0 = 1000 * 0. So if the hypothesis was 'a divides b', this should be True. But it's False. What about 'b divides a'? 0 divides 1000? No (division by zero). So if the condition is 'b divides a', then f(1000, 0) would be False. Let's re{-}check previous examples with 'b divides a': f({-}10, 10): 10 divides {-}10? Yes. f(2, {-}1): {-}1 divides 2? Yes. But f(2, {-}1) is False. So 'b divides a' is not it. Let's re{-}check 'a divides b': f({-}10, 10): {-}10 divides 10? Yes. f(2, {-}1): 2 divides {-}1? No. f({-}10, 1000): {-}10 divides 1000? Yes. f({-}2, {-}99): {-}2 divides {-}99? No.} \\
            }
        }
    }
    & & \\ \\

    \theutterance \stepcounter{utterance}  
    & & & \multicolumn{4}{p{0.6\linewidth}}{
        \cellcolor[rgb]{0.9,0.9,0.9}{
            \makecell[{{p{\linewidth}}}]{
                \texttt{\tiny{[P1$\langle$GM]}}
                \texttt{Game Violation: ParseError: Response is not valid JSON.} \\
\texttt{Output ONLY:} \\
\texttt{One JSON with confidence\_score, current\_mode, input\_rationale, and hypothesis keys ...} \\
\texttt{OR} \\
\texttt{SOLVE: ```python ... ```} \\
            }
        }
    }
     \\ \\

    \theutterance \stepcounter{utterance}  
    & \multicolumn{4}{p{0.6\linewidth}}{
        \cellcolor[rgb]{0.561,0.737,0.561}{
            \makecell[{{p{\linewidth}}}]{
                \texttt{\tiny{[P1$\rangle$GM]}}
                \texttt{\{} \\
\texttt{  "current\_mode": "probing",} \\
\texttt{  "input\_rationale": "I have six examples. The divisibility hypothesis (a divides b or b divides a) seems strong but let's check the new case: f(1000, 0) = False. 1000 divides 0? Yes, 0 = 1000 * 0. So if the hypothesis was 'a divides b', this should be True. But it's False. What about 'b divides a'? 0 divides 1000? No (division by zero). So if the condition is 'b divides a', then f(1000, 0) would be False. Let's re{-}check previous examples with 'b divides a': f({-}10, 10): 10 divides {-}10? Yes. f(2, {-}1): {-}1 divides 2? Yes. But f(2, {-}1) is False. So 'b divides a' is not it. Let's re{-}check 'a divides b': f({-}10, 10): {-}10 divides 10? Yes. f(2, {-}1): 2 divides {-}1? No. f({-}10, 1000): {-}10 divides 1000? Yes. f({-}2, {-}99): {-}2 divides {-}99? No.",} \\
            }
        }
    }
    & & \\ \\

    \theutterance \stepcounter{utterance}  
    & & & \multicolumn{4}{p{0.6\linewidth}}{
        \cellcolor[rgb]{0.9,0.9,0.9}{
            \makecell[{{p{\linewidth}}}]{
                \texttt{\tiny{[P1$\langle$GM]}}
                \texttt{Game Violation: ParseError: Response is not valid JSON.} \\
\texttt{Output ONLY:} \\
\texttt{One JSON with confidence\_score, current\_mode, input\_rationale, and hypothesis keys ...} \\
\texttt{OR} \\
\texttt{SOLVE: ```python ... ```} \\
            }
        }
    }
     \\ \\

    \theutterance \stepcounter{utterance}  
    & \multicolumn{4}{p{0.6\linewidth}}{
        \cellcolor[rgb]{0.561,0.737,0.561}{
            \makecell[{{p{\linewidth}}}]{
                \texttt{\tiny{[P1$\rangle$GM]}}
                \texttt{\{} \\
\texttt{  "current\_mode": "probing",} \\
\texttt{  "input\_rationale": "I have six examples. The divisibility hypothesis (a divides b or b divides a) seems strong but let's check the new case: f(1000, 0) = False. 1000 divides 0? Yes, 0 = 1000 * 0. So if the hypothesis was 'a divides b', this should be True. But it's False. What about 'b divides a'? 0 divides 1000? No (division by zero). So if the condition is 'b divides a', then f(1000, 0) would be False. Let's re{-}check previous examples with 'b divides a': f({-}10, 10): 10 divides {-}10? Yes. f(2, {-}1): {-}1 divides 2? Yes. But f(2, {-}1) is False. So 'b divides a' is not it. Let's re{-}check 'a divides b': f({-}10, 10): {-}10 divides 10? Yes. f(2, {-}1): 2 divides {-}1? No. f({-}10, 1000): {-}10 divides 1000? Yes. f({-}2, {-}99): {-}2 divides {-}99? No.",} \\
            }
        }
    }
    & & \\ \\

    \theutterance \stepcounter{utterance}  
    & & & \multicolumn{4}{p{0.6\linewidth}}{
        \cellcolor[rgb]{0.9,0.9,0.9}{
            \makecell[{{p{\linewidth}}}]{
                \texttt{\tiny{[P1$\langle$GM]}}
                \texttt{Game Violation: ParseError: Response is not valid JSON.} \\
\texttt{Output ONLY:} \\
\texttt{One JSON with confidence\_score, current\_mode, input\_rationale, and hypothesis keys ...} \\
\texttt{OR} \\
\texttt{SOLVE: ```python ... ```} \\
            }
        }
    }
     \\ \\

    \theutterance \stepcounter{utterance}  
    & \multicolumn{4}{p{0.6\linewidth}}{
        \cellcolor[rgb]{0.561,0.737,0.561}{
            \makecell[{{p{\linewidth}}}]{
                \texttt{\tiny{[P1$\rangle$GM]}}
                \texttt{\{} \\
\texttt{  "current\_mode": "probing",} \\
\texttt{  "input\_rationale": "I have six examples. The divisibility hypothesis (a divides b or b divides a) seems strong but let's check the new case: f(1000, 0) = False. 1000 divides 0? Yes, 0 = 1000 * 0. So if the hypothesis was 'a divides b', this should be True. But it's False. What about 'b divides a'? 0 divides 1000? No (division by zero). So if the condition is 'b divides a', then f(1000, 0) would be False. Let's re{-}check previous examples with 'b divides a': f({-}10, 10): 10 divides {-}10? Yes. f(2, {-}1): {-}1 divides 2? Yes. But f(2, {-}1) is False. So 'b divides a' is not it. Let's re{-}check 'a divides b': f({-}10, 10): {-}10 divides 10? Yes. f(2, {-}1): 2 divides {-}1? No. f({-}10, 1000): {-}10 divides 1000? Yes. f({-}2, {-}99): {-}2 divides {-}99? No.",} \\
            }
        }
    }
    & & \\ \\

    \theutterance \stepcounter{utterance}  
    & & & \multicolumn{4}{p{0.6\linewidth}}{
        \cellcolor[rgb]{0.9,0.9,0.9}{
            \makecell[{{p{\linewidth}}}]{
                \texttt{\tiny{[P1$\langle$GM]}}
                \texttt{Game Violation: ParseError: Response is not valid JSON.} \\
\texttt{Output ONLY:} \\
\texttt{One JSON with confidence\_score, current\_mode, input\_rationale, and hypothesis keys ...} \\
\texttt{OR} \\
\texttt{SOLVE: ```python ... ```} \\
            }
        }
    }
     \\ \\

    \theutterance \stepcounter{utterance}  
    & \multicolumn{4}{p{0.6\linewidth}}{
        \cellcolor[rgb]{0.561,0.737,0.561}{
            \makecell[{{p{\linewidth}}}]{
                \texttt{\tiny{[P1$\rangle$GM]}}
                \texttt{\{} \\
\texttt{  "current\_mode": "probing",} \\
\texttt{  "input\_rationale": "I have six examples. The divisibility hypothesis (a divides b or b divides a) seems strong but let's check the new case: f(1000, 0) = False. 1000 divides 0? Yes, 0 = 1000 * 0. So if the hypothesis was 'a divides b', this should be True. But it's False. What about 'b divides a'? 0 divides 1000? No (division by zero). So if the condition is 'b divides a', then f(1000, 0) would be False. Let's re{-}check previous examples with 'b divides a': f({-}10, 10): 10 divides {-}10? Yes. f(2, {-}1): {-}1 divides 2? Yes. But f(2, {-}1) is False. So 'b divides a' is not it. Let's re{-}check 'a divides b': f({-}10, 10): {-}10 divides 10? Yes. f(2, {-}1): 2 divides {-}1? No. f({-}10, 1000): {-}10 divides 1000? Yes. f({-}2, {-}99): {-}2 divides {-}99? No.",} \\
            }
        }
    }
    & & \\ \\

    \theutterance \stepcounter{utterance}  
    & & & \multicolumn{4}{p{0.6\linewidth}}{
        \cellcolor[rgb]{0.9,0.9,0.9}{
            \makecell[{{p{\linewidth}}}]{
                \texttt{\tiny{[P1$\langle$GM]}}
                \texttt{Game Violation: ParseError: Response is not valid JSON.} \\
\texttt{Output ONLY:} \\
\texttt{One JSON with confidence\_score, current\_mode, input\_rationale, and hypothesis keys ...} \\
\texttt{OR} \\
\texttt{SOLVE: ```python ... ```} \\
            }
        }
    }
     \\ \\

    \theutterance \stepcounter{utterance}  
    & \multicolumn{4}{p{0.6\linewidth}}{
        \cellcolor[rgb]{0.561,0.737,0.561}{
            \makecell[{{p{\linewidth}}}]{
                \texttt{\tiny{[P1$\rangle$GM]}}
                \texttt{\{} \\
\texttt{  "current\_mode": "probing",} \\
\texttt{  "input\_rationale": "I have six examples. The divisibility hypothesis (a divides b or b divides a) seems strong but let's check the new case: f(1000, 0) = False. 1000 divides 0? Yes, 0 = 1000 * 0. So if the hypothesis was 'a divides b', this should be True. But it's False. What about 'b divides a'? 0 divides 1000? No (division by zero). So if the condition is 'b divides a', then f(1000, 0) would be False. Let's re{-}check previous examples with 'b divides a': f({-}10, 10): 10 divides {-}10? Yes. f(2, {-}1): {-}1 divides 2? Yes. But f(2, {-}1) is False. So 'b divides a' is not it. Let's re{-}check 'a divides b': f({-}10, 10): {-}10 divides 10? Yes. f(2, {-}1): 2 divides {-}1? No. f({-}10, 1000): {-}10 divides 1000? Yes. f({-}2, {-}99): {-}2 divides {-}99? No.",} \\
            }
        }
    }
    & & \\ \\

    \theutterance \stepcounter{utterance}  
    & & & \multicolumn{4}{p{0.6\linewidth}}{
        \cellcolor[rgb]{0.9,0.9,0.9}{
            \makecell[{{p{\linewidth}}}]{
                \texttt{\tiny{[P1$\langle$GM]}}
                \texttt{Game Violation: ParseError: Response is not valid JSON.} \\
\texttt{Output ONLY:} \\
\texttt{One JSON with confidence\_score, current\_mode, input\_rationale, and hypothesis keys ...} \\
\texttt{OR} \\
\texttt{SOLVE: ```python ... ```} \\
            }
        }
    }
     \\ \\

    \theutterance \stepcounter{utterance}  
    & \multicolumn{4}{p{0.6\linewidth}}{
        \cellcolor[rgb]{0.561,0.737,0.561}{
            \makecell[{{p{\linewidth}}}]{
                \texttt{\tiny{[P1$\rangle$GM]}}
                \texttt{\{} \\
\texttt{  "current\_mode": "probing",} \\
\texttt{  "input\_rationale": "I have six examples. The divisibility hypothesis (a divides b or b divides a) seems strong but let's check the new case: f(1000, 0) = False. 1000 divides 0? Yes, 0 = 1000 * 0. So if the hypothesis was 'a divides b', this should be True. But it's False. What about 'b divides a'? 0 divides 1000? No (division by zero). So if the condition is 'b divides a', then f(1000, 0) would be False. Let's re{-}check previous examples with 'b divides a': f({-}10, 10): 10 divides {-}10? Yes. f(2, {-}1): {-}1 divides 2? Yes. But f(2, {-}1) is False. So 'b divides a' is not it. Let's re{-}check 'a divides b': f({-}10, 10): {-}10 divides 10? Yes. f(2, {-}1): 2 divides {-}1? No. f({-}10, 1000): {-}10 divides 1000? Yes. f({-}2, {-}99): {-}2 divides {-}99? No.",} \\
            }
        }
    }
    & & \\ \\

    \theutterance \stepcounter{utterance}  
    & & & \multicolumn{2}{p{0.3\linewidth}}{
        \cellcolor[rgb]{0.9,0.9,0.9}{
            \makecell[{{p{\linewidth}}}]{
                \texttt{\tiny{[GM$|$GM]}}
                \texttt{Game Verdict: LOSS (turn limit reached)} \\
            }
        }
    }
    & & \\ \\

\end{supertabular}
}

\end{document}
